\documentclass{article}
\usepackage[T1]{fontenc}
\usepackage{lmodern}
\usepackage{graphicx}
\usepackage{amsmath,amssymb}
\usepackage[a4paper,margin=2cm]{geometry}
\usepackage[absolute,overlay]{textpos}
\setlength{\TPHorizModule}{1mm}
\setlength{\TPVertModule}{1mm}
\usepackage[indonesian]{babel}
\usepackage{hyperref}
\setlength{\emergencystretch}{2em}
% unit: Halaman 37

\begin{document}

% === Halaman 37 (python/native) ===

Secara matematis, mode CFB s-bit dapat dinyatakan sebagai:

yang dalam hal ini,

j = 1, 2, 3, …, m.

Xj = antrian bergeser, dengan X1 = IV

E = fungsi enkripsi

K = kunci

b = panjang blok enkripsi

s = panjang unit enkripsi

|| = operator penyambungan (concatenation)

MSB = Most Significant Byte

37

Enkripsi:

Cj = Pj MSBs(EK (Xj))

Xj+1 = LSBb – s(Xj) || Ci

Dekripsi:

Pj = Ci MSBs(EK (Xj))

Xj+1 = LSBb – s(Xj) || Ci

\end{document}
